\documentclass[../../../include/open-logic-section]{subfiles}

\begin{document}

\olfileid[pt]{sth}{spine}{recursion}
\olsection{Teorema(s) de recursão transfinita}

O primeiro passo é, contudo, confirmar que
\olref[valpha]{defValphas} é uma definição bem sucedida. De modo mais geral,
precisamos de provar que qualquer tentativa de definir por recursão transfinita
(transfinita) terá êxito. Esse é o objetivo deste texto.

\emph{Aviso: este material é delicado}. A lição geral,
no entanto, é bem simples: a indução transfinita mais o Esquema de Substituição
garantem a legitimidade de (várias versões da) recursão
transfinita.\footnote{Lembrete: todas as fórmulas e termos podem ter parâmetros (salvo indicação em contrário).}

\begin{defn}
	Seja $\tau(x)$ um termo; seja $f$ uma função; seja $\alpha$ um ordinal. Dizemos que $f$ é uma \emph{aproximação de ordem $\alpha$} para $\tau$ sse $\dom{f} = \alpha$ e $(\forall \beta \in \alpha)f(\beta) = \tau(\funrestrictionto{f}{\beta})$.
\end{defn}

\begin{lem}[Recursão limitada]\ollabel{transrecursionfun}
Para qualquer termo $\tau(x)$ e qualquer ordinal $\alpha$, existe uma única aproximação de ordem $\alpha$ para $\tau$.
\end{lem}
\begin{proof}
Mostraremos que, para qualquer $\gamma \leq \alpha$, existe uma única
aproximação de ordem $\gamma$. 

Começamos por estabelecer a unicidade. Sejam $g$ e $h$ (respetivamente) aproximações de ordem $\gamma$ e~$\delta$. Uma indução transfinita nos seus argumentos mostra que
$g(\beta) = h(\beta)$ para qualquer $\beta \in \dom{g} \cap \dom{h} =
\gamma \cap \delta = \min(\gamma, \delta)$. Logo as aproximações são únicas (se existem) e coincidem em todos os valores.

Para estabelecer a existência, usamos agora uma indução transfinita simples
(\olref[ordinals][opps]{simpletransrecursion}) nos ordinais
$\delta \leq \alpha$. 

A função vazia é trivialmente uma aproximação de ordem $\emptyset$. 

Se $g$ é uma aproximação de ordem $\gamma$, então $g \cup
\{\tuple{\gamma, \tau(g)}\}$ é uma aproximação de ordem $\ordsucc{\gamma}$.

Se $\gamma$ é um ordinal limite e $g_\delta$ é uma aproximação de ordem $\delta$ para todo o $\delta < \gamma$, seja $g = \bigcup_{\delta \in \gamma} g_\delta$. Isto
é uma função, pois as várias $g_\delta$ coincidem em todos os valores. E
se $\delta \in \gamma$ então $g(\delta) = g_{\ordsucc{\delta}}(\delta) =
\tau(\funrestrictionto{g_{\ordsucc{\delta}}}{\delta}) =
\tau(\funrestrictionto{g}{\delta})$.

Isto completa a prova por indução transfinita.
\end{proof}

Se nos permitirmos definir um \emph{termo} em vez de uma função,
podemos retirar a cota $\alpha$ do resultado anterior. No
enunciado e na prova do resultado seguinte, quando $\sigma$ é um termo,
definimos $\funrestrictionto{\sigma}{\alpha} = \Setabs{\tuple{\beta,
\sigma(\beta)}}{\beta \in \alpha}$.

\begin{thm}[Recursão geral]\ollabel{transrecursionschema}
Para qualquer termo $\tau(x)$, podemos definir explicitamente um termo
$\sigma(x)$ tal que
$\sigma(\alpha) = \tau(\funrestrictionto{\sigma}{\alpha})$ para qualquer
ordinal~$\alpha$.
\end{thm}

\begin{proof}
Para cada $\alpha$, por \olref{transrecursionfun} existe uma única aproximação de ordem $\alpha$, $f_\alpha$, para~$\tau$. Definimos $\sigma(\alpha)$ como $f_{\ordsucc{\alpha}}(\alpha)$. Ora:
	\begin{align*}
		\sigma(\alpha) &= 
		f_{\ordsucc{\alpha}}(\alpha) \\&= 
		\tau(\funrestrictionto{f_{\ordsucc{\alpha}}}{\alpha}) \\&= 
		\tau(\Setabs{\tuple{\beta, f_{\ordsucc{\alpha}}(\beta)}}{\beta \in \alpha}) \\&= 
		\tau(\Setabs{\tuple{\beta, f_{\ordsucc{\beta}}(\beta)}}{\beta \in \alpha})\\&=
		\tau(\funrestrictionto{\sigma}{\alpha})
	\end{align*}
notando que $f_{\ordsucc{\beta}}(\beta) = f_{\ordsucc{\alpha}}(\beta)$ para todo o $\beta < \alpha$, como no \olref{transrecursionfun}. 
\end{proof}
\noindent 
Notemos que \olref{transrecursionschema} é um \emph{esquema}. Crucialmente, não podemos
esperar que $\sigma$ defina uma função, O.s.\@, um certo tipo de
\emph{conjunto}, pois então $\dom{\sigma}$ seria o conjunto de todos os
ordinais, contradizendo o paradoxo de Burali-Forti
(\olref[ordinals][basic]{buraliforti}).

Ainda falta mostrar, contudo, que \olref{transrecursionschema}
justifica a nossa definição dos $V_\alpha$. Isto pode não ser
óbvio de imediato; mas tornar-se-á claro com uma última versão simples
da recursão transfinita.

\begin{thm}[Recursão simples]\ollabel{simplerecursionschema} 
Para quaisquer termos $\tau(x)$ e $\theta(x)$ e qualquer conjunto $A$, podemos
definir explicitamente um termo $\sigma(x)$ tal que:
\begin{align*}
	\sigma(\emptyset) &= A\\
	\sigma(\ordsucc{\alpha}) &= \tau(\sigma(\alpha)) &&
		\text{para qualquer ordinal }\alpha\\
	\sigma(\alpha) &= \theta(\ran{\funrestrictionto{\sigma}{\alpha}})&&
	\text{quando }\alpha\text{ é um ordinal limite}
\end{align*}
\end{thm}

\begin{proof}
Começamos por definir um termo, $\xi(x)$, do seguinte modo: 
%t $\xi(x) = A$ if $x$ is not a function whose domain is an ordinal;
%otherwise let:
\[
	\xi(x) = 
	\begin{cases}
			A &\text{se $x$ não é uma função cujo}\\
			&\hspace{1em}\text{domínio é um ordinal; caso contrário:}\\
			\tau(x(\alpha)) & \text{se $\dom{x} = \ordsucc{\alpha}$}\\
			\theta(\ran{x}) & \text{se $\dom{x}$ é um ordinal limite}
	\end{cases}
\]
Por \olref{transrecursionschema}, existe um termo $\sigma(x)$ tal que
$\sigma(\alpha) = \xi(\funrestrictionto{\sigma}{\alpha})$ para todo o
ordinal $\alpha$; além disso, $\funrestrictionto{\sigma}{\alpha}$ é uma
função com domínio $\alpha$. Mostramos que $\sigma$ tem as propriedades
pretendidas, por indução transfinita simples
(\olref[ordinals][opps]{simpletransrecursion}). 

Primeiro, $\sigma(\emptyset) = \xi(\emptyset) = A$. 

Segundo, $\sigma(\ordsucc{\alpha}) = \xi(\funrestrictionto{\sigma}{\ordsucc{\alpha}}) = \tau(\funrestrictionto{\sigma}{\ordsucc{\alpha}}(\alpha)) = \tau(\sigma(\alpha))$.

Por fim, $\sigma(\alpha) = \xi(\funrestrictionto{\sigma}{\alpha}) = \theta(\ran{\funrestrictionto{\sigma}{\alpha}})$, quando $\alpha$ é limite.
\end{proof}
\noindent
Ora, para justificar \olref[valpha]{defValphas}, basta tomar $A
= \emptyset$, $\tau(x) = \Pow{x}$ e $\theta(x) = \bigcup x$. Por fim, isto justifica a definição dos $V_\alpha$!{}


\end{document}
